\section{Robustness checks and protocol provenance}
\label{app:earlier-diagnostics}

This appendix retains checks that directly qualify the main results. Superseded
control fixtures and development history that support no live claim are kept
with the artifact rather than in the paper.

\subsection{Nonlinear estimator--direction robustness}

Table~\ref{tab:nonlinear-suffix-factorial} reports all four
estimator--direction combinations behind Sec.~\ref{sec:nonlinear-suffix}.

\begin{table}[H]
\centering
\scriptsize
\begin{tabular}{@{}rllrrrrr@{}}
\toprule
$\gamma$ & Estimate & Direction & ISE & $g_E\!\leq\!0$ & $|p|\!\geq\!1$ & Would clip & Beyond clean max \\
\midrule
1.15 & FO & FO & 2.632 & 0/12 & 0/12 & 0/192 & 0/192 \\
1.15 & FO & Lifted & 2.381 & 10/12 & 10/12 & 30/192 & 56/192 \\
1.15 & Lifted & FO & 3.666 & 10/12 & 10/12 & 56/192 & 96/192 \\
1.15 & Lifted & Lifted & 1.870 & 0/12 & 0/12 & 0/192 & 0/192 \\
\addlinespace
0 & FO & FO & 1.693 & 0/12 & 0/12 & 0/192 & 0/192 \\
0 & FO & Lifted & 1.659 & 0/12 & 0/12 & 0/192 & 0/192 \\
0 & Lifted & FO & 1.693 & 0/12 & 0/12 & 0/192 & 0/192 \\
0 & Lifted & Lifted & 1.662 & 0/12 & 0/12 & 0/192 & 0/192 \\
\bottomrule
\end{tabular}
\caption{Secondary nonlinear-suffix estimator--direction factorial. Counts for $g_E$ and $p$ use 12 seed-level means; the last two columns use 192 rollout conditions per cell. The crossed arms at $\gamma=1.15$ fail loop-conditioning and clean-archetype scale diagnostics, so they do not identify an estimator-only advantage. The clean maximum is a scale check, not a manifold boundary.}
\label{tab:nonlinear-suffix-factorial}
\end{table}


\subsection{IOI interaction robustness}

Table~\ref{tab:ioi-budget-appendix} reports the full nested-budget curve behind
the interaction comparison in Sec.~\ref{sec:ioi}.

\begin{table}[H]
\centering
\small
\begin{tabular}{@{}rrrrr@{}}
\toprule
Budget & Per-head & Count & $P\!\times\!E$ & All pairs \\
\midrule
20  & 0.846 & 1.579 & 1.333 & 0.795 \\
40  & 0.675 & 0.732 & 0.544 & 0.559 \\
80  & 0.496 & 0.482 & 0.406 & 0.412 \\
160 & 0.486 & 0.473 & 0.394 & 0.395 \\
\bottomrule
\end{tabular}
\caption{Fresh held-out IOI mean-effect MAE by nested measurement budget.}
\label{tab:ioi-budget-appendix}
\end{table}

\subsection{Protocol provenance}

\begin{itemize}
  \item \textbf{IOI target sensitivity.} The first measurement copy was set
  aside before inspection because its source seal omitted the head-mapping file.
  After expanding the seal, we independently recomputed the frozen predictions
  and actions and repeated the 148-mask measurement. The two copies are
  identical at the byte level; only the second enters the analysis. Under the frozen
  raw-score rule, the transformed-mean control selects negative predicted
  losses in 44, 44, and 16 pools across the three targets, while direct loss
  does so in 4, 0, and 0. We retain those scores rather than clip them after
  observing the outcomes.
  \item \textbf{Qwen Copy-v2.} Copy-v1 stopped at 92.97\% candidate accuracy,
  below its frozen 95\% gate, before attention or intervention outcomes were
  computed. Copy-v2 used fresh tokens and prompts under the same gate before
  generating the locked results in Section~\ref{sec:qwen-induction}.
\end{itemize}

\subsection{Authorization tail-loss selection}

Selecting the activation observer by tail loss chooses the same layer and ridge
setting as selecting it by mean loss. Its tail loss differs from action-only by
0.039, with paired 95\% interval $[-1.013,1.156]$. We therefore retain the mean
activation gain but do not claim a tail-loss gain over action-only.

\subsection{Qwen action loss by target}

The main Qwen result gives equal weight to three targets.
Table~\ref{tab:qwen-per-target} shows the individual comparisons. Direct loss
beats transformed mean and no-op at every target. It beats natural mean at the
first two targets and ties it within uncertainty at the largest, so the
aggregate result is not a win in every cell.

\begin{table}[H]
\centering
\scriptsize
\setlength{\tabcolsep}{4pt}
\begin{tabular}{@{}lrrrr@{}}
\toprule
Target & Direct loss & Gain vs. natural mean & Gain vs. transformed mean & Gain vs. no action \\
\midrule
$0.25s$ & 0.609 & +12.0\% [0.034, 0.141] & +14.8\% [0.035, 0.188] & +29.4\% [0.205, 0.300] \\
$0.50s$ & 0.842 & +17.6\% [0.108, 0.256] & +18.7\% [0.117, 0.283] & +51.3\% [0.824, 0.942] \\
$0.75s$ & 1.158 & \textbf{$-$0.7\% [-0.038, 0.021]} & +6.3\% [0.022, 0.146] & +55.3\% [1.331, 1.530] \\
Equal-target mean & 0.870 & +8.9\% [0.053, 0.118] & +12.6\% [0.083, 0.173] & +49.6\% [0.805, 0.908] \\
\bottomrule
\end{tabular}
\caption{Qwen Copy-v2 action loss by target, where $s=3.454$ is the confirmed
full-panel mean effect. Each gain cell gives the relative loss change and, in
brackets, a 95\% four-family-prompt-cluster-by-pool bootstrap interval for the
absolute comparator-minus-direct difference. The negative natural-mean point
estimate at $0.75s$ is retained.}
\label{tab:qwen-per-target}
\end{table}
